\chapter{Программная реализация}
\label{ch:implementation}

\section{Технологический стек}

Реализация выполнена на TypeScript --- типизированном надмножестве JavaScript.
Табл.~\ref{tab:stack} суммирует используемые технологии.

\begin{table}[H]
\centering
\caption{Технологический стек Fun-Walk}
\label{tab:stack}
\begin{tabular}{@{}lll@{}}
\toprule
\textbf{Слой} & \textbf{Технология} & \textbf{Версия} \\
\midrule
Runtime & Node.js & $\geq 18$ \\
Backend framework & NestJS & 10.x \\
Frontend framework & React & 18.x \\
Сборка frontend & Vite & 6.x \\
Стили & Tailwind CSS & 3.x \\
Карты & Leaflet, react-leaflet & 1.9 / 4.2 \\
Валидация & class-validator & 0.14 \\
Маршрутизация & OSRM (HTTP) & --- \\
\bottomrule
\end{tabular}
\end{table}

\section{Реализация математического ядра}

\subsection{Модуль geo.util.ts}

Функция \texttt{haversineKm} непосредственно кодирует формулы
\eqref{eq:haversine-a}--\eqref{eq:haversine-d}:

\begin{lstlisting}[language=TypeScript,caption={Реализация формулы Haversine},label={lst:haversine}]
export function haversineKm(a: LatLng, b: LatLng): number {
  const dLat = toRadians(b.lat - a.lat);
  const dLng = toRadians(b.lng - a.lng);
  const lat1 = toRadians(a.lat);
  const lat2 = toRadians(b.lat);

  const h =
    Math.sin(dLat / 2) ** 2 +
    Math.cos(lat1) * Math.cos(lat2) * Math.sin(dLng / 2) ** 2;

  return EARTH_RADIUS_KM * 2 * Math.atan2(Math.sqrt(h), Math.sqrt(1 - h));
}
\end{lstlisting}

Функция \texttt{projectOntoSegment} реализует проекцию \eqref{eq:projection}:
вычисляется параметр $t$, затем Haversine до точки проекции. При вырожденном
отрезке ($S = E$) возвращается расстояние до старта.

\subsection{Модуль poi-scoring.util.ts}

Функция \texttt{scorePois} объединяет формулы \eqref{eq:dot-score} и
\eqref{eq:poi-score}:

\begin{lstlisting}[language=TypeScript,caption={Оценка POI},label={lst:scoring}]
const w = buildPreferenceVector(preferences);
const wNorm = Math.sqrt(w.reduce((s, v) => s + v * v, 0)) || 1;

const dotScore = dotProduct(w, f) / wNorm;
const { perpendicularKm } = projectOntoSegment(poi.location, start, end);
const totalScore = dotScore - corridorAlpha * perpendicularKm;
\end{lstlisting}

Константа \texttt{corridorAlpha = 0.15} соответствует $\alpha$ в модели.
Функция \texttt{selectCandidatePois} фильтрует по \texttt{minScore = 0.05}
и ограничивает \texttt{maxCount = 8}.

\subsection{Построение графа}

Метод \texttt{RoutePlannerService.buildWeightedGraph()}:
\begin{enumerate}
  \item Создаёт вершины \texttt{start}, \texttt{end} и POI с полями
        \texttt{attractiveness};
  \item Вычисляет матрицу расстояний $d_H$ для всех пар;
  \item Добавляет рёбра k-NN ($k = 4$) и рёбра с $d_H \leq 3{,}5$~км;
  \item Присваивает веса по \eqref{eq:edge-weight} с \texttt{ATTRACTIVENESS\_ETA = 2.5}.
\end{enumerate}

\subsection{Алгоритм Dijkstra}

Файл \texttt{dijkstra.ts} реализует классический алгоритм с массивом-очередью.
Ключевой фрагмент релаксации:

\begin{lstlisting}[language=TypeScript,caption={Релаксация ребра в Dijkstra},label={lst:dijkstra-relax}]
const alt = (dist.get(u) ?? Infinity) + weight;
if (alt < (dist.get(v) ?? Infinity)) {
  dist.set(v, alt);
  previous.set(v, u);
  if (!queue.includes(v)) queue.push(v);
}
\end{lstlisting}

Восстановление пути выполняется обратным обходом \texttt{previous} от \texttt{end}
к \texttt{start}. Функция \texttt{countEdges} делит сумму степеней на 2 для
неориентированного графа.

\section{Интеграция с OSRM}

Сервис \texttt{OsrmRoutingService} инкапсулирует HTTP-взаимодействие:
\begin{itemize}
  \item формирование URL с координатами \texttt{lng,lat} (порядок OSRM!);
  \item парсинг GeoJSON geometry из ответа;
  \item извлечение \texttt{distance} (м) и \texttt{duration} (с);
  \item метод \texttt{routeFootWithFallback} с сегментной склейкой;
  \item \texttt{deduplicateGeometry} с $\varepsilon = 10^{-5}$.
\end{itemize}

При недоступности OSRM \texttt{RoutePlannerService} возвращает
\texttt{routingSource: 'direct'} и waypoints из вершин графа --- пользователь
видит предупреждение в UI.

\section{Реализация REST-слоя}

\texttt{RoutesController} обрабатывает POST \texttt{/api/routes/plan}:

\begin{lstlisting}[language=TypeScript,caption={Фрагмент контроллера},label={lst:controller}]
@Post('routes/plan')
planRoute(@Body() dto: PlanRouteDto): PlannedRoute {
  return this.routesService.planRoute(dto);
}
\end{lstlisting}

\texttt{RoutesService} генерирует UUID, автоматическое имя маршрута
(два топ-предпочтения: «Прогулка: парки + чистый воздух») и сохраняет объект
в \texttt{Map}.

Bootstrap (\texttt{main.ts}) включает глобальный \texttt{ValidationPipe}
(\texttt{whitelist: true}, \texttt{transform: true}) и CORS.

\section{Реализация frontend}

\subsection{API-клиент}

Модуль \texttt{api/client.ts} содержит типизированные функции
\texttt{planRoute}, \texttt{fetchPois}, \texttt{fetchRoutes} и др.
Базовый путь \texttt{/api} проксируется Vite.

\subsection{Компонент MapView}

Использует \texttt{Polyline} react-leaflet для отображения \texttt{waypoints}.
При новом маршруте вызывается \texttt{fitBounds} для автоматического масштаба.
POI отображаются маркерами с popup (название, тип, AQI, шум).

\subsection{Компонент RouteStats}

Визуализирует \texttt{RouteMetrics}: дистанция, время, зелёное покрытие,
средний AQI, итоговый score (кольцевой progress). Блок \texttt{highlights}
перечисляет посещённые POI.

\subsection{Управление состоянием}

Состояние хранится в \texttt{useState} корневого \texttt{App.tsx}:
\texttt{preferences}, \texttt{start/end}, \texttt{currentRoute}, \texttt{savedRoutes}.
Отдельный state manager не используется --- достаточно для масштаба прототипа.

\section{Организация исходного кода}

Структура репозитория:
\begin{verbatim}
Fun-Walk/
  backend/src/routes/planner/   # алгоритмы
  backend/src/routes/data/      # POI Москвы
  frontend/src/components/      # UI
  frontend/src/api/             # HTTP-клиент
  scripts/                      # dev.ps1, dev.sh
  coursework/                   # пояснительная записка (LaTeX)
\end{verbatim}

Сборка production: \texttt{npm run build} (backend \texttt{nest build}, frontend
\texttt{vite build}).

\section{Выводы по главе}

Программная реализация непосредственно следует математической модели главы~\ref{ch:math}:
формулы закодированы в модулях \texttt{math/}, Dijkstra --- в \texttt{graph/},
интеграция OSRM --- в \texttt{routing/}. TypeScript обеспечивает типобезопасность
на всех уровнях. Frontend реализует интуитивный UI без дублирования бизнес-логики.
